\honest{Artificial Addition}{Eliezer Yudkowsky}{2007}

Suppose people had evolved the ability to add, and had no idea how it worked. Philosophers write ``Plus-Of(Seven, Six) = Thirteen.'' Calculators are lookup tables typed in by experts, and fail on ``two hundred plus two hundred.'' Philosophers say calculators only simulate addition, and experts say it will be at least a hundred years before calculators add as well as a twelve-year-old.

Then I list what the dissenters say. Store enough facts and real addition will appear; no, the machine must learn arithmetic as children do; no, it must read it off the Web. Build it by evolution; make it emerge from the bottom up; wait until Moore's law gives calculators the power of a brain, ``in the year 2031 on April 27 between 4:00 and 4:30 in the morning.'' Scan a brain; train neural networks without knowing how they work. G\"odel's theorem shows the brain must use quantum physics; Searle's ``Chinese Calculator Experiment'' shows there is no real addition in the system. \nb{Several voices echo real positions: Lenat's Cyc project of hand-entered common sense, Lucas and Penrose's argument from G\"odel's theorem, Searle's Chinese Room. Later, language models, neural networks trained on large amounts of text, did learn to add, and in 2025 Kantamneni and Tegmark worked out how three of them do it.}

My morals: beware assertions you cannot regenerate from your own knowledge, and beware dancing around basic confusions. Lest anyone accuse me of generalizing from fictional evidence, I say both lessons can be drawn from the real history of AI. \nb{The reader of the parable already knows that arithmetic is a simple formal system, so every voice looks foolish by construction. Whether intelligence has a key of that kind is what is in question.} The calculators play back ``knowledge'' they cannot generate, like someone repeating ``Light is waves.'' ``More on this theme tomorrow.'' \nb{That post, ``Truly Part of You,'' appears in the first part of this book.}

Whether you call the gap ``emergent'' or ``unknowable,'' you are avoiding the admission that an insight is missing. The only way to get one is years of work, and ``It's not a pursuit that academia is set up to permit.'' \nb{The one real insight the post goes on to describe, Judea Pearl's graphical models, came from academia. Pearl was at the University of California, Los Angeles, and Pearl's 1982 paper on the method acknowledges support from the National Science Foundation.} Pearl's graphical models explained what dozens of non-monotonic logics had patched case by case.

Clever ideas look promising only because you cannot see the obstacles. ``Until you know your idea will work, it won't.'' \nb{In the parable, human arithmetic was produced by evolution, the very method one of the mocked voices proposes, and nothing knew in advance that it would work.} I derive this lesson ``From the history of previous key insights in Artificial Intelligence.'' \nb{The post describes one such insight.}
